\section{Introduction}
\label{sec:intro}

Quantum error correction protects quantum information by encoding a logical qubit into many physical qubits and repeatedly measuring a set of commuting stabilizer operators~\cite{gottesman1997,nc2010}. A classical decoder processes the measurement outcomes and infers a Pauli correction. Topological codes are attractive because their stabilizers act on geometrically local qubits, and their logical error rate falls exponentially with the code distance once the physical error rate is below a threshold~\cite{kitaev2003,dennis2002,fowler2012}. Stabilizer codes are now operated on superconducting, trapped-ion and neutral-atom processors~\cite{ryananderson2021,postler2022,google2023,bluvstein2024}, and the surface code has been operated below threshold~\cite{google2025}.

Every fault-tolerant computation starts with the preparation of logical states. For a CSS code such as the surface code, the state $\ket{\bar{0}}$ is prepared by initializing every qubit in $\ket{0}$ and measuring the stabilizers for $d$ rounds~\cite{dennis2002,fowler2012}. The product state already satisfies the $Z$-type stabilizers, so their first outcomes are deterministic and act as detectors. Codes whose stabilizers mix Pauli types do not have this structure in the computational basis. A general alternative is to measure every stabilizer and the target logical operator, decode the outcomes and apply a Pauli correction. We refer to this approach as measurement, decoding and recovery (MDR). Mid-circuit measurement and feedforward of this kind were used to prepare a toric code ground state in constant depth on a trapped-ion processor~\cite{iqbal2024}. Whether an MDR protocol is fault tolerant depends on the initial state, on the circuit that extracts the syndrome and on the decoder.

In this paper we study MDR for the $XYZ^2$ hexagonal stabilizer code~\cite{Sri2022}. The code is a matching code~\cite{wootton,wootton2022} built on the honeycomb lattice of Kitaev's model~\cite{kitaev2006}. It has weight-six plaquette stabilizers, weight-two link stabilizers and weight-three boundary stabilizers, and it encodes one logical qubit into $2d^2$ physical qubits with distance $d$. The code is the concatenation of a YZZY surface code with the $[[2,1,1]]$ phase-flip codes on its links. Maximum-likelihood decoding under code-capacity noise shows low logical failure rates for $Z$-biased noise~\cite{Sri2022}. A sequential decoder that first corrects the link codes and then passes conditional error probabilities to the surface code reaches thresholds of $18.3\%$ for code-capacity depolarizing noise and $3.4\%$ for phenomenological noise~\cite{Sri2025}. Both studies assume perfect or phenomenologically noisy stabilizer measurements.

An MDR protocol is fault tolerant only if three conditions hold. The initial state must fix enough stabilizers that errors before the first measurement are detected, the extraction circuit must not spread single faults into logical errors, and the decoder must use the outcomes of all rounds together. The simplest initial state for $\plusL$, the product state $\ket{+}^{\otimes n}$, fixes only the links and $\Xb$, and it limits the circuit-level fault distance to two at every code distance. The main contributions of this paper are as follows:
\begin{enumerate}
  \item We construct a fault-tolerant MDR protocol. The data qubits start in a product state chosen in the CSS frame of the code, so that $\Xb$ and $d^2$ independent stabilizers are deterministic. All $2d^2-1$ generators are measured with a depth-six schedule selected by an exhaustive search, the full space-time syndrome is decoded, and the recovery is a Pauli frame update. The protocol has circuit-level fault distance $d+1$, the largest possible for $\plusL$.
  \item We introduce hierarchical correlated matching. This decoder treats the gauge checks, which include the links, as a lower level and the frame-deterministic checks as an upper level, and couples the two by correlated matching~\cite{fowler2013}. We compare it with a circuit-level version of sequential decoding~\cite{Sri2025}, with erasure passing between the two levels, with belief-matching~\cite{higgott2023} and with the search-based Tesseract decoder~\cite{beni2025}.
  \item We determine how the threshold depends on the number of measurement rounds. Under circuit-level depolarizing noise and with distances up to 19, the threshold of matching falls from $1.8\%$ with one round to $0.43\%$ when the number of rounds grows with the distance.
  \item We compute thresholds for six circuit-level noise models, including $Z$-biased noise and the EM3 model of hardware with native pair measurements, for which we give an extraction circuit with cat-state ancillas. Hierarchical correlated matching raises the threshold by $21\%$ to $55\%$ over matching, with the largest gains for biased noise, and belief-matching reaches $0.86\%$ under strong bias.
  \item We introduce the coset free-energy decoder, which compares the probabilities of the two logical classes instead of searching for the most likely error. A two-class ordered-statistics search, local moves between equivalent errors and an entropy correction make it as accurate as Tesseract at $d=5$. Under strongly biased noise it fails on $28\%$ fewer shots than belief-matching at $d=7$, and its advantage grows with the distance.
  \item We build one-parameter noise models of the Quantinuum Helios and H2 processors, in which the two-qubit error rate sets the scale and the other rates keep the ratios of the data sheets~\cite{helios_datasheet,h2_datasheet}. We check them against independent measurements~\cite{ransford2026} and the vendor's emulator~\cite{quantinuum_emulator}, show that the measurement crosstalk moves the crossing points of consecutive distances to smaller error rates as the code grows, and design $d=3$ and $d=5$ experiments for Helios.
\end{enumerate}

The paper is organized as follows. \autoref{sec:code} reviews the $XYZ^2$ code and its block structure. \autoref{sec:mdr} describes MDR and defines the deterministic checks and the fault distance. \autoref{sec:protocol} presents the fault-tolerant protocol and \autoref{sec:decoding} the decoders. \autoref{sec:results} contains the circuit-level simulations and \autoref{sec:hardware} the implementation on Quantinuum hardware. We conclude in \autoref{sec:conclusion}.
